\documentclass[11pt,a4paper]{article}
\usepackage[T1]{fontenc}
\usepackage[utf8]{inputenc}
\usepackage[czech]{babel}
\usepackage{lmodern,microtype,amsmath,amssymb,geometry,enumitem,tikz}
\geometry{top=17mm,bottom=17mm,left=22mm,right=22mm}
\pagestyle{empty}\setlength{\parindent}{0pt}\setlength{\parskip}{4pt}
\begin{document}
\begin{center}{\Large\bfseries Kvadratické funkce – souhrnné procvičování (102)}\end{center}
\textbf{Klíč pro učitele.} Výsledky a postupové poznámky; při náčrtu požaduj vrchol, osu a kontrolní body.\par
\section*{A – Základy}
\noindent\textbf{Úloha 1.} Funkce $f(x)=x^2$. Určete $f(-3)$, $f(0)$, $f(2)$ a její obor hodnot.\par
\[\begin{aligned}f(-3)=9,\ f(0)=0,\ f(2)=4,\ H(f)=\langle0;+\infty).\end{aligned}\]\par
\noindent\textbf{Úloha 2.} U $f(x)=-3x^2$ určete orientaci, vrchol, $H(f)$ a extrém.\par
\[\begin{aligned}Míří dolů,\ V=[0;0],\ H(f)=(-\infty;0\rangle,\ \max f=0.\end{aligned}\]\par
\noindent\textbf{Úloha 3.} Najděte $a$, jestliže $f(x)=ax^2$ prochází bodem $[2;12]$.\par
\[\begin{aligned}12=4a\Rightarrow a=3.\end{aligned}\]\par
\noindent\textbf{Úloha 4.} Pro $f(x)=(x-4)^2-2$ určete vrchol, osu a směr posunu od $y=x^2$.\par
\[\begin{aligned}V=[4;-2],\ x=4,\ \text{posun o 4 vpravo a 2 dolů}.\end{aligned}\]\par
\noindent\textbf{Úloha 5.} Pro $g(x)=-\frac12(x+2)^2+1$ určete vrchol, $H(g)$ a maximum.\par
\[\begin{aligned}V=[-2;1],\ H(g)=(-\infty;1\rangle,\ \max g=1.\end{aligned}\]\par
\section*{B – Vrchol, graf a průsečíky}
\noindent\textbf{Úloha 6.} Upravte $f(x)=x^2-6x+8$ na vrcholový tvar; určete vrchol a nulové body.\par
\[\begin{aligned}f(x)=(x-3)^2-1=(x-2)(x-4),\ V=[3;-1],\ K=\{2;4\}.\end{aligned}\]\par
\noindent\textbf{Úloha 7.} Upravte $g(x)=x^2+2x-8$; určete vrchol, osu a průsečíky s osami.\par
\[\begin{aligned}g(x)=(x+1)^2-9=(x+4)(x-2),\ V=[-1;-9],\ x=-1,\ P_y=[0;-8],\ x\in\{-4;2\}.\end{aligned}\]\par
\noindent\textbf{Úloha 8.} U $h(x)=-x^2+2x+3$ určete vrchol, obor hodnot a nulové body.\par
\[\begin{aligned}h(x)=-(x-1)^2+4,\ V=[1;4],\ H=(-\infty;4\rangle,\ K=\{-1;3\}.\end{aligned}\]\par
\noindent\textbf{Úloha 9.} Pro $p(x)=2x^2-8x+6$ určete vrchol, osu, $H(p)$ a nulové body.\par
\[\begin{aligned}p(x)=2(x-2)^2-2=2(x-1)(x-3),\ V=[2;-2],\ H=\langle-2;+\infty),\ K=\{1;3\}.\end{aligned}\]\par
\noindent\textbf{Úloha 10.} U funkce $q(x)=x^2+4x+5$ určete vrchol a vysvětlete počet nulových bodů.\par
\[\begin{aligned}q(x)=(x+2)^2+1,\ V=[-2;1],\ H=\langle1;+\infty),\ K=\varnothing.\end{aligned}\]\par
\noindent\textbf{Úloha 11.} Pro $r(x)=3x^2-5x+2$ určete průsečík s osou $y$ a ověřte, zda body $[1;0]$ a $[2;4]$ leží na grafu.\par
\[\begin{aligned}P_y=[0;2],\ r(1)=0\Rightarrow[1;0]\ \text{ano},\ r(2)=4\Rightarrow[2;4]\ \text{ano}.\end{aligned}\]\par
\noindent\textbf{Úloha 12.} Pro $s(x)=(x-1)^2$ s omezeným $D=\langle-1;3\rangle$ určete $H(s)$.\par
\[\begin{aligned}\min=0\text{ v }x=1,\ \max=4\text{ v }x=-1,3,\ H=\langle0;4\rangle.\end{aligned}\]\par
\noindent\textbf{Úloha 13.} V jedné soustavě načrtněte $y=x^2$, $y=2x^2$, $y=-\frac12x^2$. Vysvětlete vliv $a$.\par
\[\begin{aligned}V=[0;0]\ \text{u všech};\ 2x^2\text{ užší vzhůru},\ -\frac12x^2\text{ širší dolů}.\end{aligned}\]\par
\noindent\textbf{Úloha 14.} Do jedné soustavy zakreslete $y=(x+1)^2-2$ a $y=-(x-2)^2+1$; uveďte vrcholy.\par
\[\begin{aligned}V_1=[-1;-2],\ V_2=[2;1];\text{ obě paraboly jsou souměrné kolem svých os}.\end{aligned}\]\par
\section*{C – Výzvy}
\noindent\textbf{Úloha 15.} Parabola má vrchol $V[1;-2]$, je popsána $y=a(x-1)^2-2$ a prochází bodem $[0;1]$. Určete její předpis.\par
\[\begin{aligned}1=a(0-1)^2-2\Rightarrow a=3,\ f(x)=3(x-1)^2-2=3x^2-6x+1.\end{aligned}\]\par
\noindent\textbf{Úloha 16.} Funkce $f(x)=x^2-4x+c$ protíná osu $y$ v bodě $[0;5]$. Určete $c$, vrchol a počet nulových bodů.\par
\[\begin{aligned}c=5,\ f(x)=(x-2)^2+1,\ V=[2;1],\text{ nulové body nejsou}.\end{aligned}\]\par
\noindent\textbf{Úloha 17.} Parabola $y=a(x+2)^2+3$ prochází bodem $[-1;1]$. Určete $a$, vrchol a obor hodnot.\par
\[\begin{aligned}1=a+3\Rightarrow a=-2,\ V=[-2;3],\ H=(-\infty;3\rangle.\end{aligned}\]\par
\noindent\textbf{Úloha 18.} Najděte průsečíky grafů $f(x)=x^2-2x$ a $g(x)=x$.\par
\[\begin{aligned}x^2-2x=x\Rightarrow x(x-3)=0,\ x\in\{0;3\},\ P_1=[0;0],\ P_2=[3;3].\end{aligned}\]\par
\end{document}
